\begin{definition} 
    $\lsp = \dsp{\log n}$.
\end{definition}

\begin{theorem} 
    $\tree = \{ G\ |$ неорентированный граф $G$ является деревом$\} \in \lsp$.
\end{theorem}

\begin{proof}
    Дерево — связный граф, в котором число ребер = число вершин - 1. Второе условие легко проверить на логарифмической памяти (проход по некоторому представлению входного графа, что будет использовать $O(1)$ памяти). Проверку на связность заменим на проверку двух условий: 1) нет изолированных вершин; 2) при некотором хитром обходе, начатом с произвольного ребра, оно впервые повторится через $2(n - 1)$ шагов.

    Если дерево положить на плоскость и обходить его «по правилу левой руки», то путь пройдёт по каждому ребру дважды: в одну и в другую сторону -- прежде чем зациклиться. Таким образом, длина цикла будет равна $2e$, где $e$ -- число рёбер. При этом более длинным цикл быть не может ни для какого графа: если одно ребро повторилось вместе с направлением, то и дальнейшее будет повторяться. Таким образом, для дерева длина цикла равна $2(n - 1)$, а для других графов меньше. Проверять планарность можно по критерию Понтрягина-Куратовского, используя теорему Рейнгольда, однако мы хотим избежать сложной техники. Поэтому вместо проверки планарности мы определим обход без опоры на плоскую структуру.

    Обход:
    
    Пусть все вершины графа пронумерованы: $1, 2, ..., n$.
    Будем говорить, что за ребром $(i, j)$ следует ребро $(j, k)$, если $k = \min\{ l > i\ | $ ребро $(j,l)$ есть в графе$\}$. 
    Будем считать, что индексы упорядочены по циклу: после $n$ снова идет $1$ и т.д., (то есть по циклу $1 > n$ и тд.). В частности, если пришли в лист, то вернёмся.
    
    Утверждается, что для дерева цикл при таком обходе будет всегда иметь длину $2(n - 1)$. Это доказывается по индукции: если в дереве одно ребро, то очевидно. Иначе в дереве есть лист $i$. Пусть он соединён с вершиной $j$ и находится в списке между её соседей: $k$ и $l$ (т.е. $k$ — максимальный номер соседа $j$, меньший $i$, а $l$ - минимальный номер соседа $j$, больший $i$ (с учетом зацикливания). При этом случай $k = l$ не исключён). Для дерева с вычеркнутой вершиной $i$ по предположению индукции есть обход длины $2(n - 2)$, в котором есть цепочка $(k, j, l)$. По описанной процедуре в новом графе она должна замениться на цепочку $(k, j, i, j, l)$. Таким образом, прибавится два ребра, и общее число рёбер в пути составит $2(n - 1)$, что и требовалось.
    
    Алгоритм:
    \begin{itemize}
        \item Проверить, что число рёбер на единицу меньше числа вершин.
        \item Проверить, что нет изолированных вершин: проверить выходит ли из каждой вершины хотя бы одно ребро (особый случай: 1 вершина и 0 рёбер).
        \item Запустив обход с произвольного ребра, проверить, что оно впервые повторится через $2(n - 1)$ шагов.
    \end{itemize}
    
    Если все три проверки прошли, то граф — дерево, иначе нет. Для каждого из шагов достаточно логарифмической памяти: для первого и второго это совсем просто, для третьего это объясняется тем, что достаточно хранить только начальную, текущую и предыдущую вершины, а также общее число уже сделанных шагов.
\end{proof}